% ============================================================
%  H. Høffding, Søren Kierkegaard som Filosof
%  Kjøbenhavn: P. G. Philipsens Forlag, 1892.
%  Chapter IV, Section A: Erkendelsesteori (pp. 58--70)
%  TRANSCRIPTION (Danish)
%  Source: digitized epub (Det Kongelige Bibliotek),
%  cross-checked against PDF scan (kb.dk).
%  Note: Orthography restored to 1892 original (aa for aa,
%  capitalised nouns) from the modernised epub source.
% ============================================================
\documentclass[12pt,a4paper]{article}
\usepackage[utf8]{inputenc}
\usepackage[T1]{fontenc}
\usepackage[danish]{babel}
\usepackage{microtype}
\usepackage{setspace}
\usepackage[margin=1.2in]{geometry}
\usepackage{fancyhdr}
\usepackage{hyperref}

\hypersetup{
  pdftitle={Søren Kierkegaard som Filosof, Kap. IV.A: Erkendelsesteori},
  pdfauthor={H. Høffding},
  hidelinks
}

\pagestyle{fancy}
\fancyhf{}
\rhead{Høffding, \emph{Søren Kierkegaard som Filosof} (1892), IV.A}
\cfoot{\thepage}

\setstretch{1.3}

\title{\textbf{Søren Kierkegaard som Filosof}\\[6pt]
  \large Kapitel IV: Søren Kierkegaards Filosofi\\[4pt]
  A. Erkendelsesteori}
\author{Harald Høffding}
\date{Kjøbenhavn: P. G. Philipsens Forlag, 1892\\
  \small (pp.\ 58--70)}

\begin{document}
\maketitle
\thispagestyle{fancy}

\bigskip

\subsection*{Søren Kierkegaards Erkendelsesteori}

\medskip

\noindent\textbf{1.} Naar Kierkegaard drøfter det Spørgsmaal, hvilken
Erkendelse et existerende Menneske kan naae, tænker han ikke paa al
mulig Erkendelse, men kun paa hvad han kalder den væsentlige
Erkendelse, hvorved han mener den Erkendelse, der særligt forholder
sig til det erkendende Individ selv i dets Existensforhold, den
etisk-religiøse Erkendelse. Al anden Viden --- den empiriske, den
mathematiske og den historiske --- er ham ligegyldig, fordi den ikke
har noget med Existensen at gjøre. Allerede her gjør han en meget
vigtig Distinktion, anbringer et Enten-Eller. Thi det er dog et stort
Spørgsmaal, om ikke meget af hin Viden, han anser for ligegyldig, er
--- og har været --- af indgribende Betydning for den etisk-religiøse
Erkendelse. Naar denne ikke er den samme til alle Tider, saa skyldes
det for en Del hin tilsyneladende »ligegyldige« Viden, der bevidst og
ubevidst faar Indflydelse ogsaa paa de Regioner af Aandslivet, som
søges stængte for den. Men hin Distinktions Rigtighed antager
Kierkegaard nu engang for given, og dermed afstikker han Grænsen for
sin Undersøgelse.

\medskip

\noindent\textbf{2.} Tankegangen i Kierkegaards Erkendelsesteori er
rettet polemisk baade mod Hegel og mod dem, der (som Martensen) mente,
at naar man blot stillede sig paa Troens Grund og havde faaet sin
Bevidsthed genfødt, saa havde man derved faaet Betingelserne for en
højere, hidtil uopnaaelig Erkendelse. Overfor Hegel bemærkes, at han
egentlig slet ikke har svaret paa Kants Indvendinger mod en Erkendelse
af den absolute Virkelighed. Tænkningen kan ikke naae Virkeligheden:
thi i samme Øjeblik, som den vilde have naaet den, har den omsat den
til tænkt Virkelighed, eller til Mulighed; Tænkningen kan ikke komme
ud over sig selv. Overfor den spekulative Theologi bemærkes, at det
ikke interesserer Troen at forstaa sig selv paa anden Maade end ved
stadig at blive i Troens Lidenskab. Om det Højeste i Verden kan man
kun være vidende som lykkelig eller som ulykkelig Elsker; ved at
optræde med spekulativ Nyfigenhed faar man i Virkeligheden intet at
vide. Tag den Lidenskab bort, der stedse paa ny vækkes af Uvisheden
som en gavnlig Tugtemester, og Troen er borte med det samme. Den
Existerende er stadig i Vorden, og Tiden er aldrig afsluttet; derfor
kan fuld Vished aldrig naas, og vor Viden kan aldrig komme ud over
Tilnærmelser og Hypotheser, hvormed vi paa den uvæsentlige Erkendelses
Omraader ogsaa fuldt vel kan lade os nøje. Men skal det Evige og
Ubetingede gribes og fastholdes, kan det kun ske i Lidenskab, fordi
Afslutning ikke kan naas. Trangen til noget Færdigt og Afsluttet er af
det Onde, skyldes Magehed eller Utaalmodighed. Ud over den evige
Stræben efter Sandhed kommer vi ikke. --- Kierkegaard slutter sig med
Begejstring til Lessings berømte Ord, at om Gud i sin højre holdt al
Sandhed, og i sin venstre den stedse levende Stræben efter Sandhed,
vilde han gribe den venstre Haand, fordi den evige Sandhed ikke er for
et endeligt Væsen. Lessings Begrundelse af dette sit Valg er dog en
noget anden end Kierkegaards. Lessing gaar ud fra, at det ikke er
Besiddelsen, men kun Erhvervelsen, der giver Mennesket Værd, idet dets
Kræfter udvides ved at søge og det derved voxer i Fuldkommenhed. Det
er en mere positiv og psychologisk Begrundelse end den, vi i det
Følgende vil finde hos Kierkegaard, der er for meget Asket og for
optaget af at forberede sin Lære om Paradokset til at følge Lessing
paa hans naturlige Vej.

Nærmere udvikler Kierkegaards Erkendelsesteori sig gennem Begrundelse
af de to Sætninger: et logisk System kan der gives; men et
Tilværelsens System kan der ikke gives.

\medskip

\noindent\textbf{3.} Et logisk System kan der gives. Vi kan udvikle en
sammenhængende Række af almindelige Bestemmelser, som gjælder for vor
Tænken. Men her maa man da vel vogte sig for at indsmugle Begreber,
der er hentede fra Virkeligheden. Det logiske System er en Hypothese,
hvis Gyldighed for Virkeligheden er problematisk. Og det maa særligt
undersøges, om vi danner denne Hypothese efter at have lært
Virkeligheden at kjende gjennem Erfaring, altsaa som en Slags
forkortet Udtryk (Abbreviatur) for denne, eller om vi danner den
uafhængig af Erfaring. --- Kierkegaard udtaler Ønsket om at behandle
dette sidste Spørgsmaal udførligere ved en anden Lejlighed; men det er
ikke sket. Det er, som man let seer, Spørgsmaalet, om den empiristiske
eller den rationalistiske Opfattelse af Tænkningens Grundforudsætninger
er den rette.

En særegen Vanskelighed ved Opstillingen af det logiske System ligger
for Kierkegaard i, hvorledes man kan komme til at begynde paa det,
d.v.s.\ faa fat i et Udgangspunkt. Thi de simpleste Forudsætninger
opdager vi jo ved Reflexion, ved at gaa tilbage fra de afledede
Antagelser til Antagelser, hvorpaa de grunde sig --- men naar vi er
gaaet tilstrækkelig langt tilbage, hvorledes kan vi standse
Reflexionen? Denne har jo ligesaa lidt Grund til at standse paa det
ene Punkt, som paa det andet, og vil altsaa, overladt til sig selv,
blive ved i det Uendelige. Eller, som det ogsaa kan udtrykkes: Tvivlen
kan ikke standse sig selv. Der maa ske et Brud, en vilkaarlig
Afbrydelse. Der maa fattes en Beslutning om at standse den
tilbagegaaende Tankerække og fastslaa det Punkt, man er naaet til, som
Princip og Udgangspunkt for Systemet. Men hvor bliver saa
Forudsætningsløsheden af? Thi Begyndelsen paa det logiske System
betinges jo saa ved noget Andet end det Logiske, nemlig ved det, der
faar mig til at fatte Beslutningen!

»Hvad om vi saa i Stedet for at tale eller drømme om en absolut
Begyndelse, tale om et Spring?«

--- Overfor Hegel, der vilde give et absolut Tankesystem, der afsluttede
sig i sig selv som en Cirkel, er denne Kierkegaards Tankegang
rammende. Hvis man sammenligner den med den overfladiske
Katekismustale, med hvilken Martensen »gik ud over Hegel«, har man
Forskjellen mellem dem som Tænkere given i et enkelt Exempel. --- Dog
er det et Spørgsmaal, om her er saa stor Vanskelighed ved
»Begyndelsen«, som Kierkegaard mener. Vi finder ganske rigtig vore
sidste Forudsætninger ved Reflexion eller som Descartes udtrykte sig,
ved Tvivl, eller som man vist allerbedst udtrykker sig, ved Analyse.
Men i vor Søgen efter Forudsætninger har vi stedse en bestemt Opgave
for Øje, som skal løses ved de søgte Forudsætningers Hjælp. Opgaven er
at forstaa den i Erfaring givne Tilværelse eller en Del af den. Vi har
derfor en bestemt Maalestok for, hvor langt vi behøve at gaa --- nemlig
indtil vi har fundet hvad der er nødvendigt til Løsningen af hin
Opgave. Kan vi gaa videre endnu, opdager vi fjernere liggende
Forudsætninger, der frembyder sig med samme logiske Nødvendighed som
de først fundne; men der vil da foreløbig ikke være nogen Brug for
disse Forudsætninger. Maaskee fremskridende Erfaring kan føre til at
gjøre Brug af dem. --- Det er altsaa ingen Vilkaarlighed, der danner
Logikkens Begyndelse, og dens første Grundsætninger har mere end rent
arbitrær Gyldighed.

\medskip

\noindent\textbf{4.} Et Tilværelsens System kan ikke gives. Det kommer
for det første af, at da al Existens er i Tiden, i Vorden. Særligt det
sidste maa ikke glemmes. Den spekulative Filosofi, som har ment at
forklare Alt, har glemt, at den, der danner det spekulative System,
selv om han er nok saa ubetydelig (»et saadant lille bitte Dingeldangel
som den existerende Hr.\ Professor, der skriver Systemet«), dog ogsaa
hører med til den Tilværelse, der skal forklares, og han er jo aldrig
færdig med sin Existens. Systemet vilde først blive muligt, naar man
kunde se tilbage paa Existensen som afsluttet --- men det vilde
forudsætte, at man ikke længere existerede! Som vi allerede har hørt:
Livet leves forlæns, men forstaaes baglæns. --- Og selv den baglæns
Forstaael se, Indsigten i Nødvendigheden af Fortidens Udvikling turde
bero paa en Illusion: Thi hvorledes kan det, der saa længe det er
fremtidigt, kun er muligt, blive nødvendigt ved at høre Fortiden til?
Skulde det Forbigangne være mere nødvendigt end det Tilkommende? Vi
forstaar jo dog først Fortiden ret ved at gjøre os samtidige med den
--- men saa ser vi den jo som vordende, ikke afsluttet, altsaa endnu
for en stor Del blot som mulig! (Se herom »Philosophiske Smuler«.)

For den Existerende er desuden Tænken og Væren ikke ét, men sondrede.
Vi tænker enten tilbage paa det, som har været, eller frem paa det,
som vil komme. Samtidighed, absolut Enhed af Tænken og Væren er umulig
for et existerende Væsen, og existerer kun som Fantasteri. Der kan ikke
komme Kontinuitet i vor Tænken, da denne stadig afbrydes af Opgaaen i
Existensen. Al Enhed og Kontinuitet er kun Abstraktion.

Og endelig: kun det Enkelte, Konkrete existerer; at existere betyder at
være Enkelt. Men kun det Almene og Abstrakte er Tankens Gjenstand.
Ogsaa heri ligger et Misforhold mellem Tænken og Existens.

Fordi et Tilværelsens System ikke kan dannes af vor Tænken, følger dog
nu efter Kierkegaard ikke, at et saadant System ikke existerer. Der er
givet et Tilværelsens System --- men kun for »den, som selv er udenfor
Tilværelse og dog i Tilværelsen, --- som i sin Evighed er for evigt
afsluttet og dog indeslutter Tilværelsen i sig; det er Gud.« (Da
Kierkegaard definerer at existere ved at vorde, siger han konsekvent:
Gud existerer ikke, da han er evig.) --- Men for Mennesket er
Tilværelsen irrational, er et Paradox. Paradoxet, det Urimelige og
Modsigende, er efter Kierkegaard ingen Indrømmelse, men er, som han
udtrykker sig, en »ontologisk Bestemmelse«, betegner Forholdet mellem
en existerende erkendende Aand og den evige Sandhed.

Det maa indrømmes, at der er etableret et Paradox ved den skarpe
Modsigelse mellem den evige Sandhed (Tilværelsens evige Afsluttethed i
Gud) og den i stadig Vorden stedte Existens. Thi hvorledes kan fra
noget som helst Synspunkt og for nogen Tanke, om hvilken vi kan danne
os et Begreb, det staa som afsluttet, der stadig vorder? Kierkegaard
hævder bestemt, at »Bevægelse lader sig ikke tænke sub specie æterni«
[d.\ e.\ fra Evighedens Synspunkt, bortset fra Tiden]; men hvorledes
kan da Gud være vidende om Existensen, om den hele i Vorden stedte
Tilværelse, da jo »Existens ikke lader sig tænke uden Bevægelse«? ---
Der er her paa voldsom Maade sat to Tanker sammen, hvilke Kierkegaard
begge fordrer fastholdte. Erkendelsesteorien maa sikkert her vende hans
Yndlingsbestemmelse Enten-Eller imod ham og sige: enten maa Begrebet om
det absolute System i Gud være en Illusion, eller ogsaa maa den stadige
Vorden være et Skin. --- Kierkegaard behøvede, hvis han vilde vælge det
første Alternativ, derfor ikke at opgive sine theistiske Forudsætninger.
Han kunde som andre philosophiske Theister tage den Konsekvens, at Gud
selv maa være i Vorden, at Guds Væsen ikke er afsluttet. Denne Idee
udvikledes samtidig med Kierkegaards philosophiske Produktion af den
moderne philosophiske Theismes egentlige Grundlægger C.\ H.\ Weisse.
Derved vil den skarpe Modsætning være ophævet --- men det absolut
Afsluttede i Gudsideen vilde ganske vist være ofret med det samme. ---
Dersom man holder sig til, at den Tilværelse, vi kjende, udfolder sig i
Tiden, og at vi ikke kan gjøre os nogensomhelst Forestilling om nogen
anden Form for Tilværelse, saa bliver der ikke Tale om noget afsluttet
System. Men Tilværelsen bliver saa heller intet Paradox. Den bliver
irrational i mathematisk Betydning, idet kun en approximativ Erkendelse
af den bliver mulig. Den staar for os som $\sqrt{2}$, hvilken vi stedse
kan bestemme med flere Decimaler uden at udtømme den. Tilværelsen kunde
jo være saa rig og mangesided, at der var Stof nok i den for den mest
energiske og omfattende Tanke, vi kan forestille os. Men det er noget
helt andet, end at den skulde være paradoxal, det vil sige, at den
skulde staa for os med uovervindelige Modsigelser. Tilværelsens
Uudtømmelighed for Tanken svarer nøjagtigt til Lessings Idee om den
evige Stræben efter Sandhed. Naar vi fastholde denne Idee, maa vi som
objektivt Modstykke antage en stadig vordende, og i sin Vorden
uoverskuelig Tilværelse.

Kierkegaard gaar ud fra, at vi ikke, ogsaa hvad den »væsentlige«
Erkendelse angaar, kan nøjes med Tilnærmelser og Hypotheser. Hvorfor
dog ikke? Naar det blot er virkelige Tilnærmelser, og naar det er
Hypotheser, der stedse mere verificeres ved Erfaring? Da har vi Grund
nok til ikke at opgive vort Haab om at vinde klarere Forstaael se og
om praktisk at kunne bestaa i Kampen med Verden. Og det er jo et
saadant Haab, som bærer al vor Erkendelsestrang. --- Men det vil vise
sig i det Følgende, at Kierkegaard behandler Begrebet Tilnærmelse paa
en mærkelig Maade.

\medskip

\noindent\textbf{5.} At den »væsentlige« Sandhed er paradoxal, begrundes
hos Kierkegaard ved, at Tanken om Gud som den absolut evige skal
fastholdes midt i Existensens uafbrudte Vorden. Derpaa beror det
ogsaa, at den væsentlige Sandhed kun kan fastholdes med Troens
Lidenskab.

Spørger man, hvorledes Gudsideen kommer frem for Kierkegaard, da er
hans Svar, at Gud er et Postulat, et Nødværge, uden hvilket vi ikke kan
udholde at leve. Et Billede, som godt vilde udtrykke hans Tanke, vilde
være følgende: Gud er et Nødanker, som vi behøve paa Existensens vilde
Hav, og i hvilket der under Bølgegangen stadig rykkes. Vi behøve et
absolut fast Holdepunkt, en Krog, paa hvilken Tilværelsens Kjæde kan
hæftes, et Punkt, hvor den uendelige Reflexion kan finde Hvile. Men
hvad der postuleres, er kvalitativt forskjelligt fra det postulerende
Væsen: end ikke Begrebet Existens skal jo passe paa Gud! Gud forholder
sig efter Kierkegaard ikke som et superlativisk Ideal til Mennesket.
»Mellem Gud og Menneske er der en absolut Forshjel«; »Kvaliteterne er
absolut forskjellige«. --- Men saa gaar det da saaledes, at det
Nødværge, Mennesket i Existensen angst tyr til, selv gjør nye
Vanskeligheder og forøger Lidelsen: idet det nu gjælder at forene de
absolut forskjellige Kvaliteter, at tænke dem sammen! --- Troen paa Gud
opstaar ved, at Mennesket »uendelig interesseret spørger efter en anden
Virkelighed end sin egen« --- men derved sættes da strax nye Problemer.

Man forstaar, at Kierkegaard efter sin Opfattelse ikke kunde have saa
let som Martensen ved at bygge spekulativ Theologi paa Troens Grund.
Denne Grund var af vulkansk Natur, var i stadig Bevægelse, ligesom den
selv skyldte sin Opstaae n til en Revolution. Her var da ikke Tid og
Sted til at bygge, men kun til bekymret at kjæmpe for at bevare hvad
der var.

En objektiv Vished, en Erfaringsbekræftelse kan aldrig faas. Ikke blot
er Fremtiden, som ofte bemærket, stedse lige uvis, men naar man
betragter Naturen for at finde Bekræftelse paa Guds Almagt og Visdom,
finder man ved Siden af Spor i denne Retning ogsaa Spor i modsat
Retning. Af den objektive Uvished kommer vi altsaa ikke ud --- og da
det dog gjælder Liv og Død, omfattes denne objektive Uvished med
uendelig Lidenskab. Sandheden kan kun fastholdes i Subjektivitetens
Form, som den personlige Følelses og Lidenskabs Gjenstand.

\medskip

\noindent\textbf{6.} Efter at Kierkegaard saaledes har udviklet, at
Sandheden kun kan gribes og fastholdes paa subjektiv, personlig Maade,
--- at Sandheden er Subjektiviteten, vender han Sætningen om og hævder:
Subjektiviteten er Sandheden. Thi kun det, som gribes med subjektiv
Energi og Lidenskab, kan være Sandheden. Dersom der er to Mennesker,
af hvilke den ene beder i personlig Usandhed til den sande Gud, den
anden »med Uendelighedens hele Lidenskab« beder til en Afgud, saa beder
hin i Virkeligheden til en Afgud, medens denne i Sandhed beder til Gud.
Det kommer først og fremmest an paa et Hvorledes, ikke paa et Hvad, paa
den personlige Inderlighed og Energi, ikke paa det Indhold, til hvilket
denne Inderlighed knytter sig. Og med den objektive Uvished, med den
Modsigelse, Indholdet frembyder for den, der i personlig Existens skal
fastholde det, stiger Inderligheden og Lidenskaben. Sandheden er et
Vovestykke. Til Sætningen, at Subjektiviteten er Sandheden, svarer den
anden Sætning, at Sandheden (objektivt set) er Paradoxet.

Ved denne sin berømte Sætning, at Subjektiviteten er Sandheden, har
Kierkegaard skærpet den Lessingske Sætning, at kun i evig Stræben kan
man forholde sig til Sandheden. Hele Værdien flyttes over fra den
objektive (dogmatiske) til den subjektive (psychologiske) Pol. Meningen
er egentlig den: kun da har Sandheden Værdi, naar den tilegnes og
fastholdes i personlig Følelse og Stræben. Det er Sandhedens personlige
Næringsværdi, der er det Afgørende. Maalestokken ligger i den
Bevægelse, der vækkes i Personlighedens indre Liv.

Egentlig har Søren Kierkegaard hermed paa sin Vis udtalt det samme,
som Ludwig Feuerbach faa Aar forud hævdede i »Wesen des Christenthums«
i den Sætning, at Theologi er Psychologi. (\emph{Wesen des Christenthums}
udkom 1841, \emph{Uvidenskabelig Efterskrift} 1846. Kierkegaard
henviser i denne Bog, og allerede i »Stadier paa Livets Vej«,
udtrykkeligt til Feuerbachs Opfattelse af Kristendommen.) Kun ved en
Inkonsekvens kan Kierkegaard fordre et bestemt Indhold, en bestemt
objektiv Tro af Subjektiviteten: Hovedsagen er, at denne er i Tilstand
af Sandhedsstræben; hvad den forholder sig til, maa være ligegyldigt.
Det maa være det personlige Livs Fylde og Højde, som det kommer an paa,
og det maa foreløbig staa som et aabent Spørgsmaal, om et vist bestemt
dogmatisk Trosindhold er nødvendigt for et saadant Livs Udvikling.

Men han tager ikke den fulde Konsekvens af sin Sætning. Det vilde have
ført ham til at proklamere den frie Personligheds Princip. Han opgiver
ikke Antagelsen af en fra Personligheden selv forskjellig Maalestok for
Sandheden, som for ham er given i den kirkelige Dogmatik. Og dog ---
ved et nærmere Eftersyn (som vi vil foretage, naar vi har lært hans
Etik at kjende) vil det vise sig, at denne dogmatiske Rettesnor ikke er
den sidste. Han er for meget Tænker til at kunne slaa sig til Ro med
den. Hans sidste Maalestok er af formal Natur: Forholdet til det
dogmatisk-kirkelige Trosindhold er for ham den højeste Form for
personligt Liv, fordi der ved dette Forhold kommer den stærkeste
Modsigelse, den største Lidelse og Lidenskab frem i Subjektiviteten.
Den sidste og afgørende Anke mod ham vil da vise sig at være den, at
han har løsrevet Personligheden fra den reale, naturlige og sociale
Sammenhæng, i hvilken dens Liv ene kan faa et positivt Indhold. Thi
kun i den frie, med Virkeligheden kjæmpende og af Virkelighed lærende
Subjektivitet kan Sandheden bo, --- ikke i den Subjektivitet, der
gjennem formelle, asketiske Øvelser piner sig ind i Lydighed mod et
naturstridigt, livsfjendtligt Indhold. --- Men kun ved Undersøgelsen af
Kierkegaards Etik kan dette Punkt klarere belyses.

\bigskip

Det store ved Søren Kierkegaards Erkendelsesteori er den Alvor, med
hvilken han opfatter Tænkningens Sammenhæng med det personlige Liv. Der
er noget Græsk hos ham. Han vil tilbage fra den abstrakte Tænken og
Kontemplation til personlig, existentiel Tænken.

Vore Ideer blive kun alt for let til Soldater i Fredstid. Søren
Kierkegaard har ført Ideerne ud fra Kontemplationens Hygge til
Livskampens Strid og Uro og paavist den Arena, hvor den »væsentlige«
Erkendelses Gyldighed skal staa sin Prøve. Havde den dogmatiske Hemsko
ikke været, vilde hans Tanker have faaet mere direkte universel
Betydning; nu maa der ske en Oversættelse, en Udskillelsesproces, før
de kan komme til at spille den Rolle for Menneskeheden, som tilkommer
dem.

\end{document}
